\begin{frame}[t]{Ordem monomial e Pseudodivisão}
\small

    \begin{defbox}{Ordem Monomial}
Uma ordem monomial em $\kk[x_1,\ldots, x_n]$ é uma relação de ordem total $\preceq$ sobre o conjunto de monômios \[\mathbb{M}_n = \left\{\prod_{i=1}^nx_i^{\alpha_i} ; \alpha_1, \dots, \alpha_n \in \mathbb{N}\right\}\] de $\kk[x_1,\ldots, x_n]$, tal que, se $m_1,m_2 \in \mathbb{M}_n$ são tais que $m_1 \preceq m_2$, então $m_1m_3 \preceq m_2m_3, ~ \forall  m_3 \in \mathbb{M}_n$, com a propriedade de todo subconjunto não vazio de $\mathbb{M}_n$ admitir um menor elemento com respeito à $\preceq$.
    \end{defbox}
\pause
Com uma ordem $\preceq$ fixada, cada polinômio não nulo $f$, com conjunto de monômios $\displaystyle  \mathbb{M}(f)=\left\{\prod_{i=1}^nx_i^{\alpha_i}\right\}$, possui um monômio líder $\displaystyle  \mathrm{ml}(f)=\max_\preceq\left\{\prod_{i=1}^nx_i^{\alpha_i}\in\mathbb{M}(f)\right\}$ bem definido, com termo líder $\displaystyle\mathrm{tl}(f) = \max_\preceq \left\{x_i^{\alpha_i} \in \mathbb{M}(f)\right\}$. \\
\vspace{0.4cm}
Utilizaremos aqui a chamada \textit{ordem lexicográfica}: $y_i \preceq x_i, \quad \forall i \leq n$.
\end{frame}

    \begin{frame}[t]{Pseudodivisão}
    \vspace{0.6cm}
    Munidos de uma forma de ordenar os monômios de um dado polinômio multivariado, podemos definir, de maneira similar ao algoritmo de Euclides, uma forma de dividir polinômios em $\KK[x_1,\dots,x_n]$.
    \vspace{0.8cm}
    \pause
    \begin{teobox}{Algoritmo da Pseudodivisão}
        Fixada uma ordem monomial $\preceq$, e dados $f,g_1,\dots,g_s \in \KK[x_1,\dots,x_n]$, com $g_i \neq 0 \quad \forall i\in\{1,\dots,s\}$ existem polinômios $q_1,\dots,q_s,r \in \KK[x_1,\dots,x_n]$ tais que $\displaystyle f = \sum_{i=1}^sq_ig_i +r$, com $\mathrm{ml}(g_i) \nmid m \quad \forall m\in \mathbb{M}(r)$ e todo $i = 1,\dots,s.$ 
    \end{teobox}
\end{frame}